\documentclass[11pt,a4paper]{article}

% ── Codificación, fuentes y lenguaje ───────────────────────────────────────────
\usepackage[utf8]{inputenc}
\usepackage[T1]{fontenc}
\usepackage[default,scale=0.95]{sourcesanspro}
\renewcommand{\familydefault}{\sfdefault}
\usepackage[spanish,es-noshorthands,es-tabla]{babel}

% ── Geometría y espaciado ─────────────────────────────────────────────────────
\usepackage[top=2.2cm, bottom=2.5cm, left=2.2cm, right=2.2cm, headheight=15pt]{geometry}
\usepackage{setspace}
\setstretch{1.18}
\usepackage{parskip}

% ── Matemáticas y unidades ────────────────────────────────────────────────────
\usepackage{amsmath}
\usepackage{amssymb}
\usepackage{siunitx}

% ── Circuitos ─────────────────────────────────────────────────────────────────
\usepackage[european, straightvoltages]{circuitikz}

% ── Tablas ────────────────────────────────────────────────────────────────────
\usepackage{booktabs}
\usepackage{tabularx}
\usepackage{array}
\usepackage{longtable}
\usepackage{colortbl}
\usepackage{enumitem}

% ── Colores, cajas e iconos ───────────────────────────────────────────────────
\usepackage[dvipsnames]{xcolor}
\usepackage{tcolorbox}
\tcbuselibrary{skins, breakable}
\usepackage{fontawesome5}
\usepackage{graphicx}

% ── Listados de código (dark-mode infográfico) ────────────────────────────────
%   Estilo base: fondo oscuro con números de línea. Los bloques lstlisting se
%   envuelven automáticamente en una caja tcolorbox redondeada.
\usepackage{listings}

% ── Secciones con barra de color (infografía) ────────────────────────────────
\usepackage{titlesec}
\titleformat{\section}
  {\normalfont\LARGE\bfseries\color{azulNoche}}
  {\colorbox{cyanNeon}{\color{fondoOscuro}\thesection}}{0.7em}{}
  [\vspace{2pt}\color{cyanNeon}\rule{\linewidth}{1.8pt}]
\titlespacing*{\section}{0pt}{24pt}{10pt}
\titleformat{\subsection}
  {\normalfont\large\bfseries\color{azulNoche}}
  {\textcolor{cyanNeon}{\thesubsection}}{0.7em}{}
  [\vspace{1pt}\color{grisLinea}\rule{\linewidth}{0.6pt}]
\titlespacing*{\subsection}{0pt}{14pt}{6pt}
\titleformat{\subsubsection}
  {\normalfont\normalsize\bfseries\color{azulNoche}}
  {\textcolor{cyanNeon}{\thesubsubsection}}{0.7em}{}
\titlespacing*{\subsubsection}{0pt}{10pt}{4pt}

% ── Cabeceras y pies de página ────────────────────────────────────────────────
\usepackage{fancyhdr}
\usepackage{lastpage}
\pagestyle{fancy}
\fancyhf{}
\renewcommand{\headrulewidth}{0pt}
\renewcommand{\footrulewidth}{0pt}
\fancyfoot[C]{%
  \begin{minipage}{\linewidth}
    \vspace{2pt}
    \color{cyanNeon}\rule{\linewidth}{1.2pt}\\[2pt]
    \centering\color{grisTexto}\footnotesize
    Página \thepage\ de \pageref{LastPage}
  \end{minipage}%
}

% ── Hiperenlaces ──────────────────────────────────────────────────────────────
\usepackage[colorlinks=true, linkcolor=azulNoche, urlcolor=cyanNeon]{hyperref}
\sloppy
\emergencystretch 3em

% ── Paleta de colores — tecnológico dark-mode ──────────────────────────────────
\definecolor{fondoOscuro}{HTML}{0D1B2A}     % Dark navy — fondo banda título
\definecolor{verdeTurquesa}{HTML}{004D40}    % Verde turquesa — bandas de marca
\definecolor{azulNoche}{HTML}{1B3A6B}       % Midnight blue — secciones, marcos
\definecolor{azulMedio}{HTML}{2A5298}       % Azul medio — variante de acento
\definecolor{cyanNeon}{HTML}{00C8E0}        % Cyan eléctrico — acentos primarios
\definecolor{verdeSignal}{HTML}{00B887}     % Verde señal — fortalezas, éxito
\definecolor{naranjaVivo}{HTML}{FF7043}     % Naranja vivo — mejoras, advertencias
\definecolor{rojoAlerta}{HTML}{E53935}      % Rojo alerta — errores
\definecolor{grisPapel}{HTML}{F4F6FA}       % Fondo tarjetas claras
\definecolor{grisLinea}{HTML}{C8D0E0}       % Bordes sutiles
\definecolor{grisTexto}{HTML}{6B7A99}       % Texto secundario
\definecolor{amarilloNota}{HTML}{FFB300}    % Notas / advertencias doradas

% Colores específicos para bloques de código (dark-mode)
\definecolor{codigoFondo}{HTML}{1E2D3D}      % Fondo del bloque de código
\definecolor{codigoComentario}{HTML}{7A8FA6} % Comentarios
\definecolor{codigoCadena}{HTML}{FF9D5C}     % Cadenas / strings
\definecolor{codigoKeyword}{HTML}{00C8E0}    % Palabras clave
\definecolor{codigoNumero}{HTML}{00E5A0}     % Números

% ── Banda de título: portada infográfica ──────────────────────────────────────
%   Diseño en 2 capas: banda de color (por defecto fondo oscuro) + franja de
%   acento cyan abajo. Uso: \bandaTitulo[color]{título}{subtítulo}.
%   El icono se puede cambiar con \renewcommand{\iconoBanda}{...} (FontAwesome).
\newcommand{\iconoBanda}{microchip}
\newcommand{\bandaTitulo}[3][fondoOscuro]{%
  \begin{tcolorbox}[
    enhanced,
    colback=#1, colframe=#1,
    boxrule=0pt, arc=8pt, outer arc=8pt,
    width=\linewidth,
    top=18pt, bottom=6pt, left=14pt, right=14pt,
    borderline south={3.5pt}{0pt}{cyanNeon},
  ]
    \begin{minipage}[c]{0.07\linewidth}
      \centering
      \textcolor{cyanNeon}{\fontsize{30}{30}\selectfont\faIcon{\iconoBanda}}
    \end{minipage}%
    \hspace{8pt}%
    \begin{minipage}[c]{0.88\linewidth}
      {\fontsize{20}{24}\selectfont\bfseries\color{white}#2}\par\vspace{5pt}%
      \textcolor{cyanNeon!80!white}{\small\faIcon{angle-right}~#3}%
    \end{minipage}
    \vspace{4pt}
  \end{tcolorbox}%
  \vspace{8pt}%
}

% ── Tarjetas de datos (infografía) ────────────────────────────────────────────
\tcbset{
  tarjetaDato/.style={
    enhanced, breakable,
    arc=6pt, outer arc=6pt,
    colback=grisPapel, colframe=azulNoche,
    boxrule=1pt,
    fonttitle=\bfseries\small, coltitle=white,
    attach boxed title to top left={yshift=-3mm, xshift=6mm},
    boxed title style={
      colback=azulNoche, colframe=azulNoche,
      arc=4pt, boxrule=0pt,
    },
    drop fuzzy shadow=azulNoche!30!white,
    top=8pt, bottom=8pt, left=10pt, right=10pt,
  },
  tarjetaIcono/.style={
    enhanced, breakable,
    arc=6pt, outer arc=6pt,
    colback=white, colframe=grisLinea, boxrule=0.8pt,
    fonttitle=\bfseries\small, coltitle=azulNoche,
    borderline west={3pt}{0pt}{cyanNeon},
    drop fuzzy shadow=azulNoche!25!white,
    top=6pt, bottom=6pt, left=10pt, right=10pt,
  },
  cajaTitulo/.style={
    enhanced, breakable,
    arc=6pt, outer arc=6pt,
    colback=azulNoche!10!grisPapel, colframe=azulNoche,
    boxrule=1pt,
    fonttitle=\bfseries, coltitle=white,
    attach boxed title to top left={yshift=-3mm, xshift=6mm},
    boxed title style={
      colback=azulNoche, colframe=azulNoche,
      arc=4pt, boxrule=0pt,
    },
    drop fuzzy shadow=azulNoche!25!white,
    top=8pt, bottom=8pt, left=10pt, right=10pt,
  },
  cajaContenido/.style={
    enhanced, breakable,
    arc=5pt, outer arc=5pt,
    colback=grisPapel, colframe=grisLinea,
    boxrule=0.8pt,
    fonttitle=\bfseries\small, coltitle=azulNoche,
    top=6pt, bottom=6pt, left=10pt, right=10pt,
  },
  cajaFortaleza/.style={
    enhanced, breakable,
    arc=6pt, outer arc=6pt,
    colback=verdeSignal!8!white, colframe=verdeSignal,
    boxrule=1pt,
    borderline west={4pt}{0pt}{verdeSignal},
    fonttitle=\bfseries\small, coltitle=verdeSignal!70!black,
    drop fuzzy shadow=verdeSignal!20!white,
    top=6pt, bottom=6pt, left=10pt, right=10pt,
  },
  cajaMejora/.style={
    enhanced, breakable,
    arc=6pt, outer arc=6pt,
    colback=naranjaVivo!8!white, colframe=naranjaVivo,
    boxrule=1pt,
    borderline west={4pt}{0pt}{naranjaVivo},
    fonttitle=\bfseries\small, coltitle=naranjaVivo!80!black,
    drop fuzzy shadow=naranjaVivo!20!white,
    top=6pt, bottom=6pt, left=10pt, right=10pt,
  },
  cajaRecomendacion/.style={
    enhanced, breakable,
    arc=6pt, outer arc=6pt,
    colback=cyanNeon!6!white, colframe=cyanNeon!60!azulNoche,
    boxrule=1pt,
    borderline west={4pt}{0pt}{cyanNeon},
    fonttitle=\bfseries\small, coltitle=azulNoche,
    drop fuzzy shadow=azulNoche!20!white,
    top=6pt, bottom=6pt, left=10pt, right=10pt,
  },
}

% ── Ícono + texto (encabezados de sección y elementos) ────────────────────────
\newcommand{\iconotexto}[2]{\textcolor{cyanNeon}{\faIcon{#1}}~#2}

% ── Listados de código: estilo dark + auto-envoltura ─────────────────────────
\lstdefinestyle{estiloCodigo}{
    backgroundcolor=\color{codigoFondo},
    basicstyle=\ttfamily\small\color{white},
    commentstyle=\color{codigoComentario}\itshape,
    stringstyle=\color{codigoCadena},
    keywordstyle=\color{codigoKeyword}\bfseries,
    numberstyle=\tiny\color{codigoComentario},
    numbers=left,
    numbersep=10pt,
    showstringspaces=false,
    breaklines=true,
    breakatwhitespace=false,
    tabsize=4,
    frame=none,
    captionpos=b,
    aboveskip=6pt,
    belowskip=4pt,
    xleftmargin=14pt,
    framexleftmargin=14pt,
    literate=
      {á}{{\'a}}1 {é}{{\'e}}1 {í}{{\'i}}1 {ó}{{\'o}}1 {ú}{{\'u}}1
      {Á}{{\'A}}1 {É}{{\'E}}1 {Í}{{\'I}}1 {Ó}{{\'O}}1 {Ú}{{\'U}}1
      {ñ}{{\~n}}1 {Ñ}{{\~N}}1 {ü}{{\"u}}1 {Ü}{{\"U}}1
      {¿}{{?`}}1 {¡}{{!`}}1
}
\lstdefinestyle{estiloPython}{
    style=estiloCodigo,
    language=Python,
}
\lstset{style=estiloCodigo}
\BeforeBeginEnvironment{lstlisting}{%
  \begin{tcolorbox}[
    enhanced, breakable,
    arc=6pt, outer arc=6pt,
    colback=codigoFondo, colframe=azulNoche,
    boxrule=1pt,
    drop fuzzy shadow=azulNoche!25!white,
    top=2pt, bottom=2pt, left=0pt, right=0pt,
  ]%
}
\AfterEndEnvironment{lstlisting}{\end{tcolorbox}}

% ── Cajas de contenido temáticas ──────────────────────────────────────────────
\newtcolorbox{cajaRecuerda}{
    enhanced, breakable,
    arc=6pt, outer arc=6pt,
    colback=azulNoche!10!grisPapel,
    colframe=azulNoche, boxrule=1pt,
    borderline west={4pt}{0pt}{cyanNeon},
    fonttitle=\bfseries\small\color{white},
    title={\faIcon{info-circle}~Recuerda},
    attach boxed title to top left={yshift=-3mm, xshift=6mm},
    boxed title style={colback=azulNoche, arc=4pt, boxrule=0pt},
    drop fuzzy shadow=azulNoche!25!white,
    left=10pt, right=10pt, top=8pt, bottom=8pt,
}

\newtcolorbox{cajaConcepto}{
    enhanced, breakable,
    arc=6pt, outer arc=6pt,
    colback=verdeSignal!8!white,
    colframe=verdeSignal, boxrule=1pt,
    borderline west={4pt}{0pt}{verdeSignal},
    fonttitle=\bfseries\small\color{white},
    title={\faIcon{lightbulb}~Concepto clave},
    attach boxed title to top left={yshift=-3mm, xshift=6mm},
    boxed title style={colback=verdeSignal!80!black, arc=4pt, boxrule=0pt},
    drop fuzzy shadow=verdeSignal!20!white,
    left=10pt, right=10pt, top=8pt, bottom=8pt,
}

\newtcolorbox{cajaEjemplo}{
    enhanced, breakable,
    arc=6pt, outer arc=6pt,
    colback=cyanNeon!5!white,
    colframe=cyanNeon!70!azulNoche, boxrule=1pt,
    borderline west={4pt}{0pt}{cyanNeon},
    fonttitle=\bfseries\small\color{fondoOscuro},
    title={\faIcon{code}~Ejemplo},
    attach boxed title to top left={yshift=-3mm, xshift=6mm},
    boxed title style={colback=cyanNeon, arc=4pt, boxrule=0pt},
    drop fuzzy shadow=azulNoche!20!white,
    left=10pt, right=10pt, top=8pt, bottom=8pt,
}

\newtcolorbox{cajaEjercicio}{
    enhanced, breakable,
    arc=6pt, outer arc=6pt,
    colback=azulMedio!7!white,
    colframe=azulMedio, boxrule=1pt,
    borderline west={4pt}{0pt}{azulMedio},
    fonttitle=\bfseries\small\color{white},
    title={\faIcon{pencil-alt}~Ejercicio},
    attach boxed title to top left={yshift=-3mm, xshift=6mm},
    boxed title style={colback=azulMedio, arc=4pt, boxrule=0pt},
    drop fuzzy shadow=azulNoche!20!white,
    left=10pt, right=10pt, top=8pt, bottom=8pt,
}

\newtcolorbox{cajaReto}{
    enhanced, breakable,
    arc=6pt, outer arc=6pt,
    colback=naranjaVivo!6!white,
    colframe=naranjaVivo, boxrule=1pt,
    borderline west={4pt}{0pt}{naranjaVivo},
    fonttitle=\bfseries\small\color{white},
    title={\faIcon{fire}~Reto},
    attach boxed title to top left={yshift=-3mm, xshift=6mm},
    boxed title style={colback=naranjaVivo, arc=4pt, boxrule=0pt},
    drop fuzzy shadow=naranjaVivo!20!white,
    left=10pt, right=10pt, top=8pt, bottom=8pt,
}

% ── Indicadores de dificultad ─────────────────────────────────────────────────
\newcommand{\facil}{\textcolor{verdeSignal}{\faIcon{circle}\,\textbf{Fácil}}}
\newcommand{\intermedio}{\textcolor{naranjaVivo}{\faIcon{circle}\,\textbf{Intermedio}}}
\newcommand{\dificil}{\textcolor{rojoAlerta}{\faIcon{circle}\,\textbf{Difícil}}}

% ─────────────────────────────────────────────────────────────────────────────

% ── Cabeceras específicas del documento ──────────────────────────────
\fancyhead[L]{\small\color{azulNoche}\textbf{Manual de instalaci\'on --- Servidor MCP S\'ismico}}
\fancyhead[R]{\small 2026-09-23}

\begin{document}

\renewcommand{\iconoBanda}{book-open}
\bandaTitulo[verdeTurquesa]{Manual de instalaci\'on del servidor MCP s\'ismico}{Pasos para montar en su m\'aquina el servidor de datos s\'ismicos --- Prof.\ Rommel Contreras}

\vspace{6pt}
\begin{tcolorbox}[tarjetaIcono, title={\faIcon{list-check}~Qu\'e lograr\'a con esta gu\'ia}]
Usted terminar\'a con un \textbf{servidor MCP} corriendo en su PC que expone su cat\'alogo de eventos
s\'ismicos (y waveforms) a otros investigadores, quienes se conectan con sus propias LLMs. El
prototipo le entrega un cat\'alogo de \textbf{prueba sint\'etico determinista} (42 eventos de la zona
de los Andes Venezolanos, incluido el doblete con sufijo {\texttt{-B}}) para que valide todo el flujo
antes de conectar sus datos reales.

\medskip
\textbf{Arquitectura entregada (3 capas):}
\begin{itemize}[leftmargin=1.6em, itemsep=2pt]
  \item Directiva (SOP): \texttt{directives/servidor\_sismico.yaml}
  \item Orquestador (decisi\'on/validaci\'on): \texttt{flujo\_servidor\_sismico.py}
  \item Ejecuci\'on (servidor MCP): \texttt{execution/servidor\_sismico.py}
\end{itemize}
\end{tcolorbox}

% ══ REQUISITOS ════════════════════════════════════════════════════════
\section{\iconotexto{computer}{Requisitos de su m\'aquina}}

\begin{tcolorbox}[tarjetaDato, title={\faIcon{terminal}~Software necesario}]
\begin{itemize}[leftmargin=1.6em, itemsep=2pt]
  \item \textbf{Python 3.10 o superior} (verificar con \texttt{python --version}).
  \item El paquete \textbf{\texttt{mcp}} en su versi\'on \textbf{v1.x} (pin \texttt{mcp<2}):
\begin{lstlisting}
python3 -m pip install "mcp<2"
\end{lstlisting}
  \item Para probar de punta a punta en su m\'aquina, un cliente MCP: \textbf{opencode}
        o \textbf{Claude Desktop} (cualquier cliente compatible con MCP sirve).
  \item \textbf{Sin API de LLM en el servidor}: el servidor solo sirve datos y herramientas.
        Cada investigador que se conecte usa su propio modelo.
\end{itemize}
\end{tcolorbox}

\begin{tcolorbox}[cajaRecomendacion, title={\faIcon{info-circle}~Por qu\'e pin \texttt{mcp<2}}]
La API que usa el prototipo es \texttt{from mcp.server.fastmcp import FastMCP} (versi\'on 1.x del
paquete \texttt{mcp}). La versi\'on 2.x renombr\'o esta clase a \texttt{MCPServer} y rompe el c\'odigo.
Instalando \texttt{mcp<2} garantiza compatibilidad total con el servidor entregado.
\end{tcolorbox}

% ══ PASO 1: OBTENER EL CÓDIGO ═════════════════════════════════════════
\section{\iconotexto{download}{Paso 1 --- Obtener el prototipo}}

El prototipo vive en el repositorio \texttt{ELECTRONICA}. Para llevarlo a su m\'aquina:

\begin{enumerate}[leftmargin=1.8em, itemsep=4pt]
  \item Clonar el repositorio (o copiar la carpeta del proyecto):
\begin{lstlisting}
git clone <ubicacion-del-repositorio-ELECTRONICA>.git
cd ELECTRONICA
\end{lstlisting}
  \item Verificar que existen las tres piezas del flujo:
\begin{lstlisting}
ls directives/servidor_sismico.yaml
ls flujo_servidor_sismico.py
ls execution/servidor_sismico.py
\end{lstlisting}
\end{enumerate}

\begin{tcolorbox}[cajaContenido, title={\faIcon{bolt}~Nota: el cat\'alogo sint\'etico ya est\'a generado}]
En el repositorio se incluye \texttt{docs/SISMOLOGIA\_MCP/catalogo\_sismico.json} como ejemplo.
Si desea regenerarlo desde cero (mismo resultado, es determinista):
\begin{lstlisting}
python3 flujo_servidor_sismico.py generar
\end{lstlisting}
\end{tcolorbox}

% ══ PASO 2: PRUEBA LOCAL ══════════════════════════════════════════════
\section{\iconotexto{flask}{Paso 2 --- Prueba local (transporte \texttt{stdio})}}

Antes de exponer nada a la red, pruebe el servidor en su propia consola:

\begin{enumerate}[leftmargin=1.8em, itemsep=4pt]
  \item Lanzar el servidor contra el cat\'alogo sint\'etico:
\begin{lstlisting}
python3 flujo_servidor_sismico.py iniciar --transporte stdio
\end{lstlisting}
  \item Desde su cliente MCP (opencode) agregue el servidor como \textbf{local}:
\begin{lstlisting}[language=json]
{
  "mcp": {
    "sismica": {
      "type": "local",
      "command": ["python3", "execution/servidor_sismico.py", "--transporte", "stdio", "--catalogo", "docs/SISMOLOGIA_MCP/catalogo_sismico.json"],
      "enabled": true
    }
  }
}
\end{lstlisting}
  \item Saludar a las \textbf{herramientas} expuestas: \texttt{eventos} (filtra por magnitud/
        profundidad/lugar), \texttt{evento} (detalle por id), \texttt{waveform} (serie sint\'etica
        determinista), \texttt{estadisticas} (resumen global del cat\'alogo).
\end{enumerate}

\begin{tcolorbox}[cajaMejora, title={\faIcon{vial}~Compruebe el doblete}]
Preg\'untele a su cliente: \emph{``\texttt{muestra el evento VE0001-B y su waveform}''}. El sufijo
\texttt{-B} modela la r\'eplica del doblete; el waveform es \textbf{determinista}: el mismo id
siempre produce la misma se\~nal.
\end{tcolorbox}

% ══ PASO 3: CONECTAR DATOS REALES ═════════════════════════════════════
\section{\iconotexto{database}{Paso 3 --- Conectar sus datos reales}}

Cuando quiera dejar el sint\'etico, entregue su cat\'alogo en \textbf{JSON} (lista de objetos) o
\textbf{CSV} (con cabecera). Campos esperados por el servidor:

\begin{tcolorbox}[tarjetaDato, title={\faIcon{cubes}~Formato del cat\'alogo}]
\begin{tabularx}{\linewidth}{@{}p{4.4cm}X@{}}
\toprule
\textbf{Campo} & \textbf{Descripci\'on} \\
\midrule
\texttt{id}          & Identificador \'unico (ej. \texttt{RM001}, \texttt{VE1030-B}). \\
\texttt{fecha\_utc}  & Fecha/hora en ISO 8601 UTC. \\
\texttt{latitud}     & Coordenada del epicentro. \\
\texttt{longitud}    & Coordenada del epicentro. \\
\texttt{magnitud}    & Magnitud. \\
\texttt{profundidad\_km} & Profundidad en kil\'ometros. \\
\texttt{lugar}       & Texto de la zona (se usa para el filtro por lugar). \\
\texttt{tipo}        & \texttt{detectado}, \texttt{doblete}, etc. (opcional). \\
\bottomrule
\end{tabularx}
\end{tcolorbox}

\begin{enumerate}[leftmargin=1.8em, itemsep=4pt]
  \item Guardar su archivo, por ejemplo \texttt{datos/catalogo\_real.json}, y lanzar:
\begin{lstlisting}
python3 flujo_servidor_sismico.py iniciar --transporte stdio --catalogo datos/catalogo_real.json
\end{lstlisting}
  \item El orquestador \textbf{valida} el cat\'alogo antes de servir: que exista, que sea una
        lista no vac\'ia y que los \texttt{id} sean \'unicos. Si falla, le indica el motivo y
        aborta con c\'odigo 1 (no sirve datos corruptos).
  \item Si sus datos vienen de otro formato (SAC, MSEED, CAT), convi\'ertalos antes a una de las
        dos estructuras aceptadas. Puede validar el estado con:
\begin{lstlisting}
python3 flujo_servidor_sismico.py estado
\end{lstlisting}
\end{enumerate}

% ══ PASO 4: EXPONERLO A LOS INVESTIGADORES ═════════════════════════════
\section{\iconotexto{network-wired}{Paso 4 --- Exponer el servidor hacia otros investigadores}}

Para que sus colegas se conecten desde fuera, use el transporte \textbf{\texttt{streamable-http}}:
el servidor queda escuchando en una URL \texttt{http://127.0.0.1:8000/mcp}. Como en Venezuela la
mayor\'ia de ISPs usan \textbf{CGNAT} (sin IP p\'ublica), el camino m\'as pr\'actico y gratuito es un
\textbf{Cloudflare Tunnel}:

\begin{enumerate}[leftmargin=1.8em, itemsep=4pt]
  \item Instalar \texttt{cloudflared} (desc\'arguelo de la web oficial de Cloudflare para su SO).
  \item Lanzar el servidor MCP apuntando a la interfaz local:
\begin{lstlisting}
python3 flujo_servidor_sismico.py iniciar --transporte streamable-http --port 8000 --catalogo datos/catalogo_real.json
\end{lstlisting}
  \item En otra terminal, abrir el t\'unel hacia ese puerto:
\begin{lstlisting}
cloudflared tunnel --url http://127.0.0.1:8000
\end{lstlisting}
  \item \texttt{cloudflared} mostrar\'a una URL tipo \texttt{https://xxxx.trycloudflare.com}.
        El \textbf{endpoint MCP} de ese t\'unel es:
\begin{lstlisting}
https://xxxx.trycloudflare.com/mcp
\end{lstlisting}
\end{enumerate}

\begin{tcolorbox}[cajaRecomendacion, title={\faIcon{globe}~Acceso alternativo con IP p\'ublica}]
Si usted s\'i dispone de IP p\'ublica o un dominio, puede usar \texttt{nginx} + \textbf{Let's
Encrypt} como proxy inverso hacia \texttt{127.0.0.1:8000}; el t\'unel de Cloudflare es la v\'ia
recomendada por no requerir IP est\'atica.
\end{tcolorbox}

% ══ PASO 5: TOKEN Y SEGURIDAD ══════════════════════════════════════════
\section{\iconotexto{lock}{Paso 5 --- Restringir el acceso (token y HTTPS)}}

\begin{tcolorbox}[cajaRecomendacion, title={\faIcon{user-secret}~Proteger el endpoint}]
El servidor acepta el flag \texttt{--token}, pero la autorizaci\'on \textbf{real} del protocolo MCP
(OAuth de FastMCP v1) queda para un despliegue institucional. La pr\'actica recomendada es:
\begin{itemize}[leftmargin=1.6em, itemsep=2pt]
  \item Exigir \textbf{HTTPS} (el t\'unel lo da gratis y autom\'atico).
  \item Para grupos cerrados: compartir el \textbf{token} por canal cifrado y rotarlo peri\'odicamente.
        El intercambio del token se documenta a los colegas junto con la URL.
  \item No abrir el puerto 8000 crudo a Internet: siempre detras del t\'unel o de un proxy.
  \item Para despliegues con mayor exigencia, proteger a nivel de t\'unel/nginx (cabecera
        \texttt{Authorization}) o implementar el \texttt{token\_verifier} OAuth del SDK.
\end{itemize}
\end{tcolorbox}

% ══ PASO 6: CÓMO SE CONECTAN LOS COLEGAS ═══════════════════════════════
\section{\iconotexto{users}{Paso 6 --- C\'omo se conectan los investigadores}}

Cada colega configura su cliente MCP apuntando a la URL del t\'unel, sin tocar su m\'aquina.

\textbf{opencode} (en \texttt{opencode.json}):
\begin{lstlisting}[language=json]
{
  "mcp": {
    "sismica": {
      "type": "remote",
      "url": "https://xxxx.trycloudflare.com/mcp",
      "headers": { "Authorization": "Bearer <token-compartido>" },
      "enabled": true
    }
  }
}
\end{lstlisting}

\textbf{Claude Desktop} (\texttt{claude\_desktop\_config.json}):
\begin{lstlisting}[language=json]
{
  "mcpServers": {
    "sismica": {
      "type": "http",
      "url": "https://xxxx.trycloudflare.com/mcp",
      "headers": { "Authorization": "Bearer <token-compartido>" }
    }
  }
}
\end{lstlisting}

\begin{tcolorbox}[cajaContenido, title={\faIcon{question-circle}~Qu\'e puede hacer cada colega}]
Cada investigador usa \textbf{su propia LLM} (OpenRouter, Gemini u otra accesible desde su pa\'is)
para preguntar en lenguaje natural: \emph{``mostrar eventos de magnitud $>3$ en T\'achira''},
\emph{``graficar el waveform de VE0001-B''}, etc. El servidor solo responde datos; los modelos son
de cada cliente.
\end{tcolorbox}

% ══ VERIFICACIÓN ═══════════════════════════════════════════════════════
\section{\iconotexto{clipboard-check}{Verificaci\'on final (lista de control)}}

\begin{tcolorbox}[tarjetaDato, title={\faIcon{check-double}~Checklist antes de distribuir la URL}]
\begin{itemize}[leftmargin=1.6em, itemsep=3pt]
  \item El cat\'alogo real fue validado por el orquestador (c\'odigo 0 y \texttt{estado} sin errores).
  \item Prueba local con \texttt{stdio} responde las 4 herramientas.
  \item Con \texttt{streamable-http} + túnel, la URL \texttt{https://.../mcp} responde al handshake.
  \item HTTPS activo y token compartido por canal seguro.
  \item Si la PC se apaga, el servidor cae: para disponibilidad robusta, considere publicar el
        cat\'alogo en GitHub o Hugging Face para que cada colega apunte a su propio clon.
\end{itemize}
\end{tcolorbox}

% ══ SOLUCIÓN DE PROBLEMAS ══════════════════════════════════════════════
\section{\iconotexto{wrench}{Soluci\'on de problemas frecuentes}}

\begin{tcolorbox}[tarjetaDato, title={\faIcon{exclamation-triangle}~Errores y acciones}]
\begin{tabularx}{\linewidth}{@{}p{7.2cm}X@{}}
\toprule
\textbf{S\'intoma} & \textbf{Acci\'on} \\
\midrule
\texttt{mcp} no importa / \texttt{ModuleNotFoundError} & Instalar: \texttt{pip install "mcp<2"}. \\
\texttt{FastMCP.run() got an unexpected keyword argument 'host'} & Usar mcp v1: host/puerto van en el
constructor \texttt{FastMCP(...)}, no en \texttt{run()}. \\
Cat\'alogo rechazado (c\'odigo 1) & Revisar formato: lista JSON o CSV con cabecera, \texttt{id}
\'unicos, campos requeridos. \\
Puerto 8000 ocupado & Cambiar de puerto y repetir (\texttt{--port 8001}). \\
URL del túnel no responde & Confirmar que \texttt{cloudflared} siga corriendo y que el endpoint sea
\texttt{/mcp}. \\
Clientes con 401/403 & Verificar el token compartido y que el túnel/nginx aplique la cabecera
correctamente. \\
\bottomrule
\end{tabularx}
\end{tcolorbox}

% ══ CIERRE ═════════════════════════════════════════════════════════════
\vspace{10pt}
\begin{tcolorbox}[cajaConcepto]
\textbf{Resumen:} el prototipo est\'a listo para que usted lo clone, lo pruebe con el cat\'alogo
sint\'etico y, cuando tenga su dato real, lo conecte y comparta la URL del túnel con la comunidad.
Si necesita integrar autenticaci\'on institucional (OAuth/SSO), el siguiente paso es implementar el
\texttt{token\_verifier} de FastMCP bajo demanda.
\end{tcolorbox}

\vspace{12pt}
\begin{center}
\footnotesize\color{grisTexto}
Generado desde \texttt{ELECTRONICA} (estilo infogr\'afico propio) \textbullet\ 2026-09-23.
\end{center}

\end{document}
